\chapter{膨張宇宙の幾何学}

要点は膨張宇宙とその幾何学である。目的は、膨張を駆動するものを問う前に、空間が膨張するとはどういうことかを記述することである。ニュートン的な\(F=ma\)宇宙論は後で扱い、その後に初めて幾何学と力学が完全なフリードマン--ロバートソン--ウォーカー枠組みへと組み立てられる。

\section{膨張する直線上の共動座標}

最も単純な例から始めよう。時空ではなく一次元空間である。等間隔に銀河が並んだ無限の直線を想像する。座標\(x\)はまだ物理的な距離を測ってはいない。それは銀河にラベルを付けるものである。\(x\)を銀河の名前と考えることができ、その名前は、配置全体が伸び縮みしても銀河に付いたままである。

\(x=0\)と呼ばれる銀河から\(x=1\)と呼ばれる銀河までの物理的距離を\(a\)とする。すると任意の二つの銀河の間の物理的距離は
\begin{equation}
d=a\,\Delta x.
\end{equation}
である。座標差\(\Delta x\)は単にラベルの差にすぎない。因子\(a\)がそのラベル差を、物差しで測られる実際の距離に変換する。

次に、この直線を伸びる弾性糸として想像しよう。ラベルは銀河に付いたままであるが、隣り合う銀河の間の距離は時間とともに変化する。同じ関係は
\begin{equation}
d(t)=a(t)\,\Delta x.
\end{equation}
となる。この段階で、一様性とは単に、隣り合う銀河がどこでも等間隔に並んでいることを意味する。ただしその共通の間隔は時間とともに変化しうる。

\section{ハッブルの法則と中心の不在}

二つの固定された銀河の相対速度を求めるには、その対を固定して\(\Delta x\)が時間とともに変化しないようにする。微分すると
\begin{align}
\dot d &= \dot a\,\Delta x \\
       &= \frac{\dot a}{a}\,a\,\Delta x \\
       &= \frac{\dot a}{a}\,d.
\end{align}
が得られる。ここで
\begin{equation}
H(t)\equiv \frac{\dot a}{a},
\end{equation}
と定義すれば、
\begin{equation}
\dot d = H(t)\,d.
\end{equation}
となる。これがハッブルの法則である。

歴史的な表現である``ハッブル定数''は誤解を招く。この量は時間的にではなく空間的に一定である。ある瞬間にはどの銀河に対しても同じであるが、一般には宇宙の進化とともに変化する。ハッブルパラメータと考えた方がよい。

この法則は、より遠い銀河ほど速く後退することを述べている。それは中心を特別扱いしない。この議論のどこにも\(x=0\)やその他の点を特権化するものはない。どの観測者も同じ規則を見る。したがってこの模型における膨張宇宙は、一つの特別な場所からの爆発ではない。

赤方偏移は直ちに現れる。二つの銀河が互いに離れながら光信号を交換すると、受け取る光は赤くなる。後退速度が大きいほど赤方偏移も大きい。したがってスペクトル線は後退速度について教えてくれる。距離の測定はより難しく、ハッブルの初期の推定は粗かった。彼の最初のプロットは完全な直線ではなく、大きな不確かさをもつ散らばった点の集合であった。それでも、幾何学的原理はノイズを生き延びた。後退速度は距離に比例するのである。

\begin{classroomqa}
\audiencequestion{銀河が十分に離れていれば、ハッブルの法則によってそれらは光速より速く動くのか。}\lecturetimestamp{00:20:23}
\lecturerresponse{はい、この後退の意味ではそうなる。君の言っていることは正しい。それについては必ず後で戻ってくる。それは最も重要な事柄の一つである。}
\end{classroomqa}

\section{高次元、等方性、そして\(a(t)\)の時間依存性}

空間次元が増えても本質的なことは何も変わらない。二次元では\((x,y)\)でラベル付けされた銀河の正方形格子を、三次元では\((x,y,z)\)でラベル付けされた銀河の格子を想像する。座標は再び共動ラベルである。それらは物理的距離を直接測るのではなく、銀河とともに動く。

\(x\)方向の間隔が\(y\)方向の間隔と同じであるべきなのはなぜか。答えは観測による。もし銀河がある方向で別の方向より系統的に密集していれば、空は方向によって異なって見えるだろう。十分に大きなスケールではそうはならない。この方向に関する同一性が等方性である。

二次元では二つの銀河の間の距離は
\begin{equation}
d=a\sqrt{(\Delta x)^2+(\Delta y)^2},
\end{equation}
であり、三次元では
\begin{equation}
d=a\sqrt{(\Delta x)^2+(\Delta y)^2+(\Delta z)^2}.
\end{equation}
である。固定された銀河の対に対して座標差は一定に保たれるので、再び微分すると
\begin{equation}
\dot d=\frac{\dot a}{a}\,d.
\end{equation}
が得られる。ハッブルの法則は変わらない。

一様性とは場所から場所への同一性を意味する。等方性とは方向から方向への同一性を意味する。どちらも小さなスケールでは厳密ではない。天の川を通して見ることはその平面の外を見ることとは異なり、個々の銀河は短い距離では一様に分布していない。しかし多くの銀河について平均すれば、大きなスケールの宇宙はどこでも、どの方向でもほぼ同じに見える。

\begin{classroomqa}
\audiencequestion{速度が距離に比例し、距離が増大しているなら、固定された銀河の一対は互いに加速しながら遠ざかっていなければならないのだろうか。}\lecturetimestamp{00:40:22}
\lecturerresponse{いいえ。ハッブルの法則が述べているのは、ある瞬間において後退速度が距離に比例するということである。それ自体は、固定された一対の速度が時間とともに増大するか減少するかについては何も述べていない。それは\(a(t)\)の詳細な時間依存性に依存する。\(a(t)\propto t\)ならば、隣り合う銀河は一定の速さで離れていく。\(a(t)\propto t^2\)ならば、固定された一対の後退速度は時間とともに増大する。\(a(t)\propto \log t\)ならば、それは時間とともに減少する。法則\(\dot d=(\dot a/a)d\)は距離の関数として常に成り立つが、加速と減速は別個の問題である。}
\end{classroomqa}

長い間、宇宙論研究者たちは宇宙は減速していると予想していた。重力は、投げ上げられた物体を減速させるのと同じように、初期の膨張を遅くするはずだからである。ここでの要点はただ、ハッブルの法則だけでは\(\dot a\)や\(\ddot a\)の振る舞いは決まらないということである。

\begin{classroomqa}
\audiencequestion{スケール因子\(a\)の初期値はどのように設定するのですか。}\lecturetimestamp{00:46:18}
\lecturerresponse{その全体的な規格化は任意である。物理的に重要なのは、\(a\)の数値的な大きさそのものではなく、それがどのように変化するかである。別のスケール因子\(b=3a\)を定義すると、\(\dot b/b=\dot a/a\)となる。\(a\)を定数で再スケールすることは、慣習や単位を変えるだけであって、ハッブルパラメータを変えるものではない。}
\end{classroomqa}

\section{空間計量から時空計量へ}

ここで\(\Delta x\)を\(dx\)に置き換える。隣り合う点の間の物理的距離は
\begin{equation}
ds=a\,dx,
\end{equation}
であるから
\begin{equation}
ds^2=a^2\,dx^2.
\end{equation}
これが一次元における空間計量である。\(dx^2\)の係数は
\begin{equation}
g_{xx}=a^2.
\end{equation}
したがって、膨張する宇宙は、計量の空間部分が時間に依存すると言うことによって記述できる。

空間から時空へ移るには、相対論的間隔を用いる。一次元の空間では、
\begin{equation}
d\tau^2=dt^2-a^2(t)\,dx^2.
\end{equation}
共動観測者に対しては\(dx=0\)であるから、\(d\tau=dt\)となる。通常の宇宙時間とは、銀河とともに動く時計によって測られる固有時である。

この時点での黒板の図は、共動世界線、線素、計量の記法をまとめて示すので有用である。
\begin{figure}[H]
\centering
\includegraphics[width=0.82\textwidth]{lecture_01_figure_02.png}
\caption{膨張する一次元宇宙における共動世界線と、対応する時空間隔。}
\end{figure}

三次元の空間では、同じパターンから慎重な拡張が得られる\footnote{黒板は00:56:38--00:58:51のあたりで一部が隠れている。この三次元の形は、一次元の空間から三次元への口頭での拡張と、黒板上で見える対角パターンに従っている。}
\begin{equation}
d\tau^2=dt^2-a^2(t)\bigl(dx^2+dy^2+dz^2\bigr),
\end{equation}
対角計量
\begin{equation}
g_{\mu\nu}=\mathrm{diag}\bigl(1,-a^2,-a^2,-a^2\bigr).
\end{equation}
非対角項は消える。
\begin{equation}
g_{tx}=g_{ty}=g_{tz}=0.
\end{equation}

一様性とは、計量の係数が時間には依存してもよいが空間位置には依存しないことを意味する。等方性とは、三つの空間方向が同じように入ることを意味する。

\begin{classroomqa}
\audiencequestion{\(dx\)を距離と考えるべきではないのですか。}\lecturetimestamp{01:01:18}
\lecturerresponse{いいえ。\(dx\)は距離ではない。\(a\,dx\)が距離である。\(dx\)は単なるラベル間隔である。これを考えるのに有用なのは、ある瞬間に誰かが銀河を通る座標の網を描き、その網が銀河に従うことを要求すると想像することである。座標の間隔は構成上固定されたままだが、物理的距離は変化する。距離の単位を\(a\)ではなく\(x\)に入れることもできる。それは慣習である。重要なのは積\(a\,dx\)である。}
\end{classroomqa}

\begin{classroomqa}
\audiencequestion{\(dt^2\)の計量係数が\(1\)に等しいのは、これも単に慣習によるのですか。}\lecturetimestamp{01:03:24}
\lecturerresponse{はい。それ自体には物理はない。係数が時間だけに依存するなら、時間変数を変えてそれがちょうど\(dt^2\)になるようにできる。同様に\(c=1\)とする単位を選んでもよい。ここで慣習でないのは、係数が位置に依存しないことである。ただし時間には依存してもよい。}
\end{classroomqa}

\begin{classroomqa}
\audiencequestion{これをアインシュタインの場の方程式に入れると、エネルギー運動量テンソルについて何が得られますか。}\lecturetimestamp{01:05:59}
\lecturerresponse{それは次に来る種類の問いであるが、動力学に属するものであって、現在の議論には属さない。この計量をアインシュタイン方程式に入れて、どのようなエネルギー運動量分布がそれを解にするかを問うことができる。それは非常に重要であるが、当面の主題は動力学ではなく幾何学である。}
\end{classroomqa}

\section{地平線が隠すもの}

大規模な宇宙が一様かつ等方であるとみなされるのは、あくまで近似としてである。恒星、惑星、あるいは個々の銀河のスケールでは、決して一様ではない。あらゆる局所的な観測において厳密に等方であるわけでもない。

農家の畑は、正しいスケール依存性を与えてくれる。あるスケールでは規則的に見えるが、はるかに小さいスケールでは明らかにそうではない。畑全体より大きいスケールでは、再び一様には見えなくなる。宇宙論も同じように成り立つ。多数の銀河について平均すれば、大規模な分布はほぼ一様に見える。十分に小さいスケールではそうではない。

宇宙は、われわれが見ることのできる部分より大きく、またわれわれが決して見ることのできない部分よりも大きい。実際に探ることのできるスケールでは、局所的な変動を別にすれば、等方かつ一様である。それを超えると、直接的な観測的保証は何もない。

\begin{classroomqa}
\audiencequestion{われわれは決して直接には突き止められないだろうとおっしゃいました。それは、間接的に突き止める方法があり得るという意味でしょうか。}\lecturetimestamp{01:10:34}
\lecturerresponse{おそらく。近ごろの大方の考えでは、十分に大きなスケールでは宇宙は一様ではない。私自身の賭けも同じである。しかしそれは、私が決して勝つことも負けることもない賭けだ。基本的な理由により、われわれは現在見えるところよりずっと遠くを見ることはできないからである。すなわち、地平線の存在である。}
\end{classroomqa}

\begin{classroomqa}
\audiencequestion{見える範囲を超えた質量を探すのに、重力を利用することはできないのでしょうか。}\lecturetimestamp{01:11:53}
\lecturerresponse{できない。遠方の質量は、われわれを単に自由落下させることができるだけである。閉じた部屋に座っていても、われわれは太陽の重力場の中にあり、その周りを公転しているが、外を見なければ、太陽がそこにあると簡単に見分ける方法はない。それを明らかにするのは潮汐力だけであり、源が十分に遠ければ、その潮汐効果は検出できないほど小さくなりうる。}
\end{classroomqa}

したがって地平線は、直接的な知識への究極の障壁である。彼はさらに、ほぼ確実に、宇宙はわれわれが決して見ることのできない領域よりも体積で少なくとも千倍大きいと述べ、その主張の根拠は後回しにしている。

\section{境界はあるが端はない}

次に閉じた幾何学に移る。無限直線は開いており、非有界である。線分は、有限の広がりしか持たないという意味で有界である。有界とは境界があるという意味ではなく、有限の広がりを持つという意味である。

壁のある立方体ではうまくいかない。それは一様性を破る。壁の近くにいる者と壁から遠くにいる者は等価ではないからである。また等方性も破る。中心にいる者でさえ、方向によって異なるものを見るからである。

円や球の内部も同様に失敗する。中心にいる者はどの方向にも同じものを見るが、その対称性はただ一つの特別な点に属するだけである。中心から離れれば、壁はある方向には近く、別の方向には遠い。するとその空間は、一般的な観測者にとって一様でも等方でもない。

非常に大きなスケールでは、宇宙は一様かつ等方であるというのが作業仮説である。これが宇宙原理である。それは有用だが、改訂から守られてはいない。

\begin{classroomqa}
\audiencequestion{至る所で等方であることは、一様であることを含意するのでしょうか。}\lecturetimestamp{01:18:57}
\lecturerresponse{はい。すべての観測者がどの方向にも同じものを見るならば、それは一様性を含意する。しかし逆は成り立たない。一様性それ自体は等方性を含意しない。宇宙は一様でありながら、なお優先方向を持つことができる。たとえば、すべて同じ向きを指す矢で満たされていればそうである。}
\end{classroomqa}

正しい有界模型は、特徴のない球の表面である。その表面に閉じ込められた二次元の生物は、世界が有限であっても、あらゆる場所が局所的に他のあらゆる場所と同じであり、あらゆる方向が他のあらゆる方向と同じであると見いだす。これが、端のない境界である。

\section{円上の閉じた一次元宇宙}

代数を簡単に保つため、球の例を一次元の空間、すなわち円上の生活に還元する。一次元では、等方性とは左と右が等価であることを意味する。円はまさにその特徴を持つ。一点を選び、それを\(\theta=0\)と呼ぶ。他のすべての点は、ラジアンで測られる角度座標\(\theta\)によってラベル付けされる。原点の一方の側すぐの点は、小さな負の角度によっても、\(2\pi\)のすぐ手前の角度によっても記述できる。この周期性は幾何学の一部である。

円の住人は、外部の半径や中心ではなく、円周にアクセスできる。彼らは、同じ点に戻るまで円に沿って進むことによって、周囲の全距離を測ることができる。周囲の平面に円を描くならば、半径に似た量を導入してもよいが、それは埋め込みの描像に属するものであり、彼らが測る内在的な幾何学に属するものではない。

ここでは黒板の図が役に立つ。角度のラベルと全周長を同じ図に保つことができるからである。
\begin{figure}[H]
\centering
\includegraphics[width=0.82\textwidth]{lecture_01_figure_03.png}
\caption{全周長\(2\pi a\)の閉じた一次元宇宙の角度による記述。}
\end{figure}

便利な慣習は\footnote{01:26:40--01:29:20頃、記法は全円周そのものをスケール因子とみなす場合と、それを\(2\pi a\)と書く場合との間で一時的に揺れる。ここでは後者の慣習を用いる。それは弧長の公式に引き継がれる慣習だからである。}
\begin{equation}
L=2\pi a,
\end{equation}
ここで\(L\)は全円周である。

小さな角度差\(\Delta\theta\)に対して、物理的距離は
\begin{equation}
d=a\,\Delta\theta.
\end{equation}
である。無限小に移行すると
\begin{equation}
ds=a\,d\theta,
\end{equation}
となり、したがって
\begin{equation}
ds^2=a^2\,d\theta^2.
\end{equation}
である。これが円上の局所計量である。

\(a\)が時間に依存するなら、閉じた一次元宇宙は膨張または収縮する。固定された一対の銀河に対して、
\begin{equation}
d(t)=a(t)\,\Delta\theta,
\end{equation}
であるから、
\begin{equation}
\dot d=\dot a\,\Delta\theta
      =\frac{\dot a}{a}\,d.
\end{equation}
となる。ここでもハッブルの法則は生き残る。少なくとも、円の反対回りが曖昧になるほど離れていない一対に対しては。

時空計量は
\begin{equation}
d\tau^2=dt^2-a^2(t)\,d\theta^2.
\end{equation}
となる。\(a\)が一定なら、時空の描像は円柱である。\(a\propto t\)なら、それは円錐である。\(a(t)\)がより不規則に変化するなら、描像は花瓶のようになる。

この描像を逆向きにたどると、銀河はどんどん近づいていく。円の周りに\(360\)個の銀河があったなら、遅い時刻にはよく広がっており、より早い時期には非常に近く集まり、さらに早くにはもっと近くなる。やがてそれらはもはや別々の銀河ではなくなり、さらに早くには別々の原子ですらなくなる。しかしそれでも、空間に特別な点は存在しない。

これが、この模型においてビッグバンがどこで起こったかという問いへの最も明快な答えである。\(a(t)\to 0\)となると、すべての間隔\(a(t)\,\Delta\theta\)はゼロに向かう。宇宙全体が至る所で同時に縮む。銀河は外部空間の特権的な位置から飛び去っているのではない。幾何学そのもののスケールが変化している。その意味で答えは、至る所である。

\begin{classroomqa}
\audiencequestion{時空は曲がっていると言いました。それが円柱だったとしても曲がっているのですか。}\lecturetimestamp{01:38:04}
\lecturerresponse{いいえ。もしそれが円柱なら、曲がってはいない。それは\(a\)が時間的に一定である場合である。曲がって見えるのは、より高次元の空間に描かれているからにすぎない。それを切り開けば、引き伸ばすことなく平らに広げられる。}
\end{classroomqa}

そこから議論は円錐へ移る。\(a\propto t\)なら、時空の描像は円錐である。その先端から離れたところでは、円錐も同じ内在的な意味で平らである。すなわち、平面に開くことができる。先端は異なる。それは円錐特異点である。

次の段階は、球状宇宙の幾何学、そして\(a(t)\)がどのように成長または縮小するかを決める力学である。
